\documentclass[preview]{standalone}

\usepackage{amsmath}
\usepackage{amssymb}
\usepackage{bettelini}
\usepackage{stellar}
\usepackage{definitions}

\begin{document}

\id{measuretheory-measure}
\genpage

\section{Measures}

\begin{snippet}{lebesgue-integration-completeness-motivation}
    The normed space \(C([0,1])\), equipped with
    \[
        \lVert f\rVert_1=\int_0^1\lvert f(x)\rvert\,dx,
    \]
    is not complete. Its completion is the space \(L^1([0,1])\), constructed
    using the Lebesgue integral. This is one reason for enlarging the classical
    theory of integration.
\end{snippet}

\begin{snippet}{measure-domain-and-codomain-motivation}
    A positive measure assigns a value in \([0,+\infty]\) to each measurable
    subset of a set \(X\). Allowing \(+\infty\) is essential for unbounded
    spaces, while restricting to nonnegative values avoids undefined expressions
    such as \(+\infty-\infty\). The domain is a \(\sigma\)-algebra rather than
    the whole power set because, in general, not every subset can be assigned a
    measure with the desired properties.
\end{snippet}

\begin{snippetdefinition}{measure-definition}{Measure}
    Let \((X,\Sigma)\) be a \measurablespace. A \emph{measure} on
    \((X,\Sigma)\) is a \function
    \[
        \mu\colon\Sigma\fromto[0,+\infty]
    \]
    such that
    \begin{enumerate}
        \item \(\mu(\emptyset)=0\);
        \item for every sequence \(\{A_n\}_{n=1}^{\infty}\subseteq\Sigma\) of
        pairwise disjoint sets,
        \[
            \mu\left(\bigcup_{n=1}^{\infty}A_n\right)
            =\sum_{n=1}^{\infty}\mu(A_n).
        \]
    \end{enumerate}
    The second property is called \emph{countable additivity} or
    \emph{\(\sigma\)-additivity}.
\end{snippetdefinition}

\begin{snippet}{extended-nonnegative-arithmetic-conventions}
    For nonnegative extended-real values, the conventions are
    \[
        a+(+\infty)=+\infty,
        \qquad
        a\cdot(+\infty)=
        \begin{cases}
            +\infty & a>0,\\
            0 & a=0.
        \end{cases}
    \]
    A series of nonnegative terms is independent of the order of summation,
    which makes countable additivity compatible with the fact that a union does
    not depend on the order in which its sets are listed. Signed and complex
    measures require separate definitions: one cannot allow both \(+\infty\)
    and \(-\infty\), and countable sums of complex values must converge
    unconditionally.
\end{snippet}

\section{Properties}

\begin{snippettheorem}{measure-basic-properties-theorem}{Basic properties of a measure}
    Let \((X,\Sigma,\mu)\) be a \measurespace. Then:
    \begin{enumerate}
        \item \emph{finite additivity:} if \(A_1,\ldots,A_m\in\Sigma\) are
        pairwise disjoint, then
        \[
            \mu\left(\bigcup_{i=1}^m A_i\right)=\sum_{i=1}^m\mu(A_i);
        \]
        \item \emph{monotonicity:} if \(A,B\in\Sigma\) and \(A\subseteq B\), then
        \(\mu(A)\leq\mu(B)\);
        \item if \(A\subseteq B\) and \(\mu(A)<+\infty\), then
        \[
            \mu(B\difference A)=\mu(B)-\mu(A);
        \]
        \item \emph{countable subadditivity:} for every sequence
        \(A_n\in\Sigma\),
        \[
            \mu\left(\bigcup_{n=1}^{\infty}A_n\right)
            \leq\sum_{n=1}^{\infty}\mu(A_n);
        \]
        \item \emph{continuity from below:} if \(A_1\subseteq A_2\subseteq\cdots\), then
        \[
            \mu\left(\bigcup_{n=1}^{\infty}A_n\right)
            =\lim_{n\to\infty}\mu(A_n)=\sup_n\mu(A_n);
        \]
        \item \emph{continuity from above:} if \(A_1\supseteq A_2\supseteq\cdots\)
        and \(\mu(A_m)<+\infty\) for some \(m\), then
        \[
            \mu\left(\bigcap_{n=1}^{\infty}A_n\right)
            =\lim_{n\to\infty}\mu(A_n)=\inf_n\mu(A_n).
        \]
    \end{enumerate}
\end{snippettheorem}

\begin{snippetproof}{measure-basic-properties-theorem-proof}{measure-basic-properties-theorem}{Basic properties of a measure}
    \begin{enumerate}
        \item Extend the finite family by setting \(A_i=\emptyset\) for \(i>m\),
        and apply countable additivity.

        \item Since \(B=A\sqcup(B\difference A)\), finite additivity gives
        \[
            \mu(B)=\mu(A)+\mu(B\difference A)\geq\mu(A).
        \]

        \item The same identity can be rearranged because \(\mu(A)<+\infty\).

        \item Define
        \[
            B_1=A_1,
            \qquad
            B_n=A_n\difference\bigcup_{k=1}^{n-1}A_k
            \quad(n\geq2).
        \]
        The sets \(B_n\) are pairwise disjoint, \(B_n\subseteq A_n\), and
        \(\bigcup_nB_n=\bigcup_nA_n\). Hence
        \[
            \mu\left(\bigcup_nA_n\right)
            =\sum_n\mu(B_n)
            \leq\sum_n\mu(A_n).
        \]

        \item Set \(B_1=A_1\) and \(B_n=A_n\difference A_{n-1}\) for \(n\geq2\).
        Then the \(B_n\) are pairwise disjoint,
        \[
            A_n=\bigcup_{i=1}^nB_i,
            \qquad
            \bigcup_{n=1}^{\infty}A_n=\bigcup_{n=1}^{\infty}B_n.
        \]
        Therefore
        \[
            \mu(A_n)=\sum_{i=1}^n\mu(B_i)
            \longrightarrow
            \sum_{i=1}^{\infty}\mu(B_i)
            =\mu\left(\bigcup_{n=1}^{\infty}A_n\right).
        \]
        Monotonicity makes \(\mu(A_n)\) increasing, so its limit is its supremum.

        \item Choose \(m\) such that \(\mu(A_m)<+\infty\), and for \(n\geq m\) set
        \(B_n=A_m\difference A_n\). The sequence \(B_n\) is increasing and
        \[
            \bigcup_{n=m}^{\infty}B_n
            =A_m\difference\bigcap_{n=1}^{\infty}A_n.
        \]
        By continuity from below and the finite-measure difference formula,
        \begin{align*}
            \mu(A_m)-\mu\left(\bigcap_{n=1}^{\infty}A_n\right)
            &=\mu\left(A_m\difference\bigcap_{n=1}^{\infty}A_n\right)\\
            &=\lim_{n\to\infty}\mu(A_m\difference A_n)\\
            &=\mu(A_m)-\lim_{n\to\infty}\mu(A_n).
        \end{align*}
        Cancelling the finite quantity \(\mu(A_m)\) proves the identity. By
        monotonicity, \(\mu(A_n)\) is decreasing, so its limit is its infimum.
    \end{enumerate}
\end{snippetproof}

\begin{snippetexample}{continuity-from-above-finiteness-counterexample}{Why continuity from above needs finiteness}
    On \(\naturalnumbers\) with counting measure, let
    \[
        A_n=\{n,n+1,n+2,\ldots\}.
    \]
    Then \(A_n\downarrow\emptyset\), but \(\mu(A_n)=+\infty\) for every \(n\).
    Hence
    \[
        \mu\left(\bigcap_{n=1}^{\infty}A_n\right)=0
        \neq+\infty=\inf_n\mu(A_n).
    \]
\end{snippetexample}

\section{Measure spaces and examples}

\begin{snippetdefinition}{measure-space-definition}{Measure space}
    A \emph{measure space} is a triple \((X,\Sigma,\mu)\), where
    \((X,\Sigma)\) is a \measurablespace and \(\mu\) is a \measure on
    \((X,\Sigma)\). When the \(\sigma\)-algebra is understood, one may write
    \((X,\mu)\).
\end{snippetdefinition}

\begin{snippet}{finite-measure-normalization}
    If \(0<\mu(X)<+\infty\), then
    \[
        \nu(A)=\frac{\mu(A)}{\mu(X)}
    \]
    defines a measure with \(\nu(X)=1\). Thus every nonzero finite measure can
    be normalized to a probability measure; measurable subsets are then called
    events.
\end{snippet}

\begin{snippetexample}{zero-and-infinite-measures-example}{Degenerate measures}
    The function \(\mu(A)=0\) for every \(A\in\Sigma\) is a measure. Another
    measure is
    \[
        \nu(A)=
        \begin{cases}
            0 & A=\emptyset,\\
            +\infty & A\neq\emptyset.
        \end{cases}
    \]
    These are useful boundary cases, although most applications require sets of
    positive finite measure.
\end{snippetexample}

\begin{snippetdefinition}{counting-measure-definition}{Counting measure}
    Let \(X\) be a \set and equip it with \(\powerset(X)\). The
    \emph{counting measure} is defined by
    \[
        \mu(A)=
        \begin{cases}
            \cardinality{A} & A\text{ is finite},\\
            +\infty & A\text{ is infinite},
        \end{cases}
        \qquad A\subseteq X.
    \]
\end{snippetdefinition}

\begin{plainhtml}
    With the full power set as sigma algebra, every subset and every map from the space into a measurable space is measurable.
\end{plainhtml}

\includesnpt{dirac-measure-definition}

\end{document}
